\subsection{求解与结果}
三种情景“仅禁渔”“禁渔+污染”“禁渔+污染+入侵”的积分后期平均指标和综合得分如表 \ref{tab:q5result} 所示。我们更关心的是三类指标并不同步变化：组合赋权下，三种情景的综合得分分别为 0.662、0.423 和 0.155。与仅禁渔相比，污染情景的食源指标由 0.407 升至 0.692，但江豚和长江鲟分别降至 0.636 和 0.716；在污染基础上再叠加入侵后，入侵压力升至 0.641，综合得分继续下降。
\begin{table}[htbp]
\centering
\caption{问题五的积分后期平均指标与综合得分}
\label{tab:q5result}
\begin{tabularx}{0.95\textwidth}{>{\raggedright\arraybackslash}p{3.0cm} >{\centering\arraybackslash}p{1.6cm} >{\centering\arraybackslash}p{1.6cm} >{\centering\arraybackslash}p{1.6cm} >{\centering\arraybackslash}p{1.6cm} >{\centering\arraybackslash}p{1.6cm}}
\toprule
情景 & 食源 & 江豚 & 长江鲟 & 入侵压力 & $\mathcal{H}$ \\
\midrule
仅禁渔 & 0.407 & 1.143 & 1.135 & 0.602 & 0.662 \\
禁渔+污染 & 0.692 & 0.636 & 0.716 & 0.546 & 0.423 \\
禁渔+污染+入侵 & 0.598 & 0.624 & 0.711 & 0.641 & 0.155 \\
\bottomrule
\end{tabularx}
\end{table}
\begin{table}[htbp]
\centering
\caption{问题五的赋权结果与排序稳健性}
\label{tab:q5_weight}
\small
\setlength{\tabcolsep}{2pt}
\begin{tabularx}{0.95\textwidth}{>{\raggedright\arraybackslash}p{2.2cm} >{\raggedright\arraybackslash}p{4.3cm} >{\raggedright\arraybackslash}p{3.5cm} X}
\toprule
方案 & 权重 $(M,D,A,V)$ & 排序结果 & 说明 \\
\midrule
熵权 & $(0.144,0.337,0.351,0.168)$ & 仅禁渔 \textgreater 禁渔+污染 \textgreater 禁渔+污染+入侵 & 根据指标离散程度确定权重 \\
CRITIC & $(0.319,0.242,0.244,0.194)$ & 仅禁渔 \textgreater 禁渔+污染 \textgreater 禁渔+污染+入侵 & 同时考虑对比强度与相关性 \\
组合权重 & $(0.232,0.289,0.298,0.181)$ & 仅禁渔 \textgreater 禁渔+污染 \textgreater 禁渔+污染+入侵 & 取熵权与 CRITIC 权重的平均值 \\
等权对照 & $(0.250,0.250,0.250,0.250)$ & 仅禁渔 \textgreater 禁渔+污染 \textgreater 禁渔+污染+入侵 & 用于检查排序是否依赖某一套客观权重 \\
局部扰动 & -- & 仅禁渔 \textgreater 禁渔+污染 \textgreater 禁渔+污染+入侵 & 围绕组合权重做 2000 次 $[0.85,1.15]$ 重归一扰动，完整排序稳定率为 100\% \\
\bottomrule
\end{tabularx}
\normalsize
\end{table}
污染情景中食源指标抬升而江豚、长江鲟同步下降，反映的是当前参数组下的捕食释放机制。放到方程结构里看，污染一方面通过 $u$ 下压 $K_M$，另一方面更强地抬升江豚和长江鲟的死亡项；后者造成的捕食减少在尾窗内超过了承载损失，于是 $M$ 的平均值短时回升。这个回升只对应过程量的再分配，不意味着污染改善了系统健康。

熵权、CRITIC、等权和组合赋权得到的排序完全一致，围绕组合权重做 2000 次 $[0.85,1.15]$ 重归一扰动后，完整排序稳定率和“入侵情景最差”稳定率均为 100\%。不同赋权下比较关系都没有改变，因此污染削弱禁渔收益、入侵进一步恶化系统状态的次序判断并不依赖某一套特定赋权。

\begin{figure}[htbp]
\centering
\includegraphics[width=0.88\textwidth]{fig05_q5_scenarios.png}
\caption{问题五的情景比较结果}
\label{fig:q5}
\end{figure}

